\section{Introduction}
\label{sec:intro}

Satellite and aerial archives hold information for environmental, climate, and humanitarian applications, but their size and technical complexity often prevent their effective use. Earth observation foundation models aim to lower that barrier with general-purpose backbones that downstream tasks reuse.

The Clay project\blfootnote{$^{\ast}$Corresponding author: \href{mailto:bruno@madewithclay.org}{\texttt{bruno@madewithclay.org}}. Author contributions---Making Clay~v1.5 corpus, architecture, and training (\S\ref{sec:model}): S.R.M., D.W., M.G., and B.S.-A.N.; Evaluations, benchmarks, and applications (\S\ref{sec:comparative}--\S\ref{sec:lessons}): B.S.-A.N., D.H., S.B., K.K., M.G., L.C., and A.C.; Clay stewardship and governance: A.C., L.C., D.H., and B.S.-A.N. Apache-2.0 code and weights, CC-BY-4.0 embedding catalogs: \url{https://github.com/Clay-foundation/model}, \url{https://huggingface.co/made-with-clay/Clay}, \url{https://source.coop/clay}.} released three open Earth observation (EO) foundation model checkpoints in 2024: v0 (January 2024), v1.0 (May 2024), and v1.5 (November 2024)~\cite{clay2024model,clay2024hf}. In May 2025 the project relicensed v1.5 from OpenRAIL-M to Apache-2.0. The project never published a technical report~\cite{clay2025issue369}. This is that report: it describes the released v1.5 model and evaluates it (\S\ref{sec:model}--\S\ref{sec:lessons}). The checkpoints are distributed with public Source Cooperative GeoParquet embedding catalogs~\cite{clay2026sourcecoop}, a TorchGeo loader for these catalogs~\cite{torchgeo2026clayembeddings}, Clay backbones in TerraTorch~\cite{terratorch2026claybackbones,gomes2025terratorch}, and an ArcGIS model~\cite{esri2025claydlpk}. Major TOM also hosts v1.5 embeddings of its Core-S2L2A 249k subset~\cite{majortom2026clayv15}. Others have evaluated Clay~v1 ($768\text{-D}$ ViT-B/8) in GEO-Bench-2~\cite{simumba2026geobench2}, where it ranks first on the Core and pixel-wise subsets, and in TerraTorch, Wiratama et al.~\cite{wiratama2025indonesia}, and AlphaEarth~\cite{brown2025alphaearth}. Clay~v1.5 ($1{,}024\text{-D}$ ViT-L/8) is evaluated in OlmoEarth~\cite{herzog2026olmoearth,olmoearth2026v13}, Gong et al.~\cite{gong2026urban}, Kaushik et al.~\cite{kaushik2026flood}, and VHR-Bench~\cite{sharma2026vhrbench}. TESSERA~v2~\cite{feng2026tesserav2} also compares a Clay checkpoint but does not state which version.

Work on pretrained vision encoders~\cite{he2022mae,oquab2024dinov2} establishes that frozen backbones paired with linear probes can replace task-specific training. EO sensors span optical, thermal, hyperspectral, and synthetic aperture radar (SAR) spectra~\cite{hong2024spectralgpt,braham2025spectralearth}, cover four orders of magnitude in ground sample distance (GSD, $0.1\text{--}500\,\text{m}$)~\cite{reed2023scalemae,irvin2023usat}, and observe Earth dynamics from flood extent to seasonal crop phenology~\cite{manas2021seco,cong2022satmae}. Current EO models~\cite{zhu2026foundations,stewart2026earthembeddings} combine masked image modeling~\cite{jakubik2023prithvi,xiong2024dofa}, contrastive or self-distillation objectives~\cite{fuller2023croma,waldmann2025panopticon}, temporal sequence objectives~\cite{tseng2023presto,astruc2025anysat}, and multimodal or coordinate-aligned training~\cite{jakubik2025terramind,klemmer2025satclip}. Apart from the training objective, most EO encoders fall into two families:
\begin{enumerate}
  \item \textbf{Sentinel-anchored and fixed-schema systems.} Temporal pixel-series models such as AlphaEarth Foundations and TESSERA~\cite{feng2026tessera} summarize multi-date observations into $10\,\text{m}$ pixel vectors ($64\text{--}128\text{-D}$) distributed as annual composites. Fixed-schema encoders (AnySat, SkySense~v2~\cite{zhang2025skysensev2}, Prithvi-EO-2.0~\cite{szwarcman2026prithvi2}, TerraMind, Galileo~\cite{tseng2025galileo}, and OlmoEarth) attach dedicated stems to predefined sensor schemas. These models do best on $10\text{--}30\,\text{m}$ phenology-driven tasks. They either ship precomputed rasters without model weights (AlphaEarth) or bind inference to fixed band sets.
  \item \textbf{Physically conditioned, sensor-agnostic encoders.} Models such as Clay, Panopticon, Copernicus-FM~\cite{wang2025copernicusfm}, and FlexiMo~\cite{li2026fleximo} generate patch projections, positional encodings, or both from wavelength and GSD at run time, following Scale-MAE, USat, ChannelViT~\cite{bao2024channelvit}, and DOFA. One frozen checkpoint then embeds unseen instruments and sub-meter to kilometer imagery without retraining.
\end{enumerate}
